\chapter{Programação}
\label{cap:programacao}

\begin{edital}
Desenvolvimento em Java/JavaEE/JakartaEE/JPA/JavaScript; frameworks JUnit,
Hibernate, JSF, Primefaces, Spring, Spring Cloud, Spring Boot; mobile
(Android/iOS); low-code/no-code; análise estática (clean code, SonarQube);
padrões XML, XSLT, UDDI, REST, JSON; DevOps; Git; codificação (transacional,
analítico, mobile, API); reuso; blockchain. (Java em si tem capítulo próprio:
Capítulo \ref{cap:java}.)
\end{edital}
\peso{MÉDIO-ALTO.}

\section{Ecossistema Spring (cai muito)}

\begin{center}
\begin{tabular}{@{}p{2.6cm}p{8.6cm}@{}}
\toprule
\textbf{Framework} & \textbf{Papel} \\
\midrule
Spring & contêiner de inversão de controle (IoC) e injeção de dependência \\
Spring Boot & Spring com autoconfiguração e servidor embutido; elimina config manual e boilerplate \\
Spring Cloud & ferramentas para sistemas distribuídos/microsserviços (config, discovery, gateway) \\
Hibernate & ORM (mapeamento objeto-relacional), implementa JPA \\
JUnit & framework de testes de unidade \\
\bottomrule
\end{tabular}
\end{center}

\begin{pegadinha}
Spring Boot \textbf{não exige} config manual de servidor (embute Tomcat) e
não exige Tomcat/JBoss standalone; ``Spring só serve monólito'' é falso.
Hibernate é \textbf{ORM/persistência}, não é framework de teste. JUnit é
teste, \textbf{não} é ORM. Spring Cloud = distribuído; Spring Boot = app
individual --- a FGV troca os papéis do ecossistema numa frase só.
\end{pegadinha}

\section{Formatos de dados: XML, JSON, XSLT}

\begin{center}
\begin{tabular}{@{}p{2.6cm}p{4.6cm}p{4.6cm}@{}}
\toprule
& \textbf{XML} & \textbf{JSON} \\
\midrule
Estrutura & marcação com tags, verboso & pares chave-valor, leve \\
Uso hoje & legado, config, SOAP & \textbf{APIs web}, transporte leve \\
\bottomrule
\end{tabular}
\end{center}

\textbf{XSLT:} linguagem que \textbf{transforma XML} em outro formato (HTML,
XML, texto). \textbf{Não} transforma JSON. \textbf{REST} usa JSON
tipicamente; \textbf{SOAP} usa XML.

\begin{pegadinha}
``XSLT transforma JSON'' = falso (é XML); ``XML e JSON são
equivalentes/idênticos'' = falso; JSON não tem comentários nem tipo date
nativo.
\end{pegadinha}

\section{Mobile e low-code}

\begin{center}
\begin{tabular}{@{}p{2.6cm}p{3.4cm}p{4.4cm}@{}}
\toprule
\textbf{Abordagem} & \textbf{O que é} & \textbf{Ferramentas / linguagem} \\
\midrule
Nativo & app específico para cada SO & Android: \textbf{Kotlin/Java}; iOS: \textbf{Swift/Objective-C} \\
Multiplataforma & um código para vários SO & \textbf{Flutter (Dart)}, \textbf{React Native (JS)}, Xamarin (.NET), Ionic (web) \\
Híbrido & web dentro de contêiner nativo (WebView) & Ionic, Cordova \\
PWA & web instalável, offline & ver Capítulo \ref{cap:frontend} \\
\bottomrule
\end{tabular}
\end{center}

\textbf{Flutter = Dart} (a FGV troca por JS/Kotlin); \textbf{React Native =
JavaScript} (renderiza componentes nativos, diferente do WebView do
híbrido). Trade-off: nativo dá melhor desempenho/acesso ao hardware; cross
reduz custo e tempo (um código só).

\textbf{Low-code/no-code:} desenvolvimento visual, pouca ou nenhuma escrita
de código; acelera entrega e permite ``citizen developers'', mas limita
customização e pode gerar \textit{lock-in} na plataforma.

\begin{pegadinha}
Pares mobile que a FGV inverte: Dart$\to$Flutter, Kotlin$\to$Android,
Swift$\to$iOS. Distrator de definição decorada: pede o framework Dart e
oferece React Native (JS), Xamarin (C\#) ou Ionic --- nenhum é Dart.
\end{pegadinha}

\section{Java EE / web server-side (JSF, Primefaces)}

\textbf{JSF (JavaServer Faces):} framework de UI \textbf{baseado em
componentes}, server-side, orientado a eventos; segue o padrão \textbf{MVC}.
Ciclo de vida de requisição com fases (restore view, apply values,
validation, update model, invoke application, render response).

\textbf{Primefaces:} biblioteca de \textbf{componentes ricos} para JSF
(tabelas, gráficos, ajax pronto).

\textbf{JSP/Servlet:} camada web mais antiga do Java EE (Servlet processa a
requisição; JSP gera a view). JSF é a evolução baseada em componentes.

\begin{pegadinha}
JSF é \textbf{server-side baseado em componentes} (diferente de SPA
client-side como React); Primefaces é biblioteca de componentes \textbf{de
JSF}, não framework à parte.
\end{pegadinha}

\section{Versionamento com Git}

\textbf{Git é distribuído:} cada clone tem o histórico completo (diferente
do SVN, que é centralizado e lida mal com binários). Conceitos: commit,
branch, merge, \texttt{git pull}/\texttt{push}, \texttt{git clone}, merge
$\times$ rebase, staging area. GitLab CI usa variáveis e pipelines.

\section{Qualidade de código}

\textbf{Clean Code:} nomes claros, funções curtas, baixo acoplamento.

\textbf{SonarQube:} análise \textbf{estática} (sem executar) --- code
smells, bugs, vulnerabilidades, dívida técnica, cobertura.

DevOps/CI-CD, reuso e padrões de projeto permeiam este bloco (ver Capítulos
\ref{cap:eng-software} e \ref{cap:arquitetura}).

\section{Blockchain}

\textbf{Cadeia de blocos} encadeados por \textbf{hash do bloco anterior}
(imutável). Cada bloco: hash anterior, timestamp, dados/transações, nonce,
raiz de Merkle. \textbf{O saldo das carteiras NÃO é armazenado no bloco} ---
é derivado do histórico de transações (UTXO no Bitcoin). Descentralizado,
consenso (PoW/PoS), contratos inteligentes (Ethereum).

\begin{pegadinha}
Pergunta clássica: ``o que NÃO é armazenado no bloco'' $\to$ resposta:
\textbf{saldo das carteiras} (é calculado a partir do histórico, não
guardado diretamente).
\end{pegadinha}

\section{Python aplicado a dados (visto em prova de TI recente)}

\begin{itemize}
\item \textbf{NumPy:} \texttt{view()} compartilha buffer (\texttt{.base}
aponta para o original) $\times$ \texttt{copy()} cria cópia independente
(\texttt{.base = None}). Alterar antes do \texttt{view()} reflete no view.
Shape de vetor 1-D termina em vírgula: \texttt{(3,)}, não \texttt{(3,1)}.
\item \textbf{matplotlib.pyplot:} parâmetro \texttt{explode} destaca fatia
da pizza.
\item \textbf{PHP 8:} funções de sessão (\texttt{session\_start},
\texttt{session\_destroy}, \texttt{session\_regenerate\_id}).
\end{itemize}

\begin{pegadinha}
NumPy: a FGV inverte \texttt{view}/\texttt{copy} --- diz que \texttt{copy}
compartilha base ou que \texttt{view} tem \texttt{base None}; ou erra o
shape do array 1-D.
\end{pegadinha}

\section{Leitura ativa de código --- técnica transversal}
\label{sec:leitura-ativa-programacao}

A FGV não para na terminologia: em Programação ela mostra um trecho pequeno
de código (Python, SQL, JS, config) e pergunta a \textbf{saída} ou o
\textbf{efeito}. Quem só decorou conceito trava aqui --- o antídoto é
\textbf{rodar o código mentalmente}, linha a linha, como um interpretador,
anotando o estado das variáveis num rascunho.

\begin{conceito}[Protocolo de leitura ativa]
\begin{enumerate}
\item \textbf{Não leia por cima.} Releia a linha inteira antes de assumir o
que ela faz --- a FGV muda \emph{um} parâmetro, \emph{um} operador ou
\emph{um} índice em relação ao que ``pareceria óbvio''.
\item \textbf{Anote o estado.} Para cada variável relevante, escreva o valor
depois de cada linha executada --- não confie na memória depois da terceira
linha.
\item \textbf{Desconfie de referência $\times$ cópia.} É a armadilha mais
comum em qualquer linguagem: duas variáveis apontam para o \textbf{mesmo}
objeto/array, e mudar uma ``muda'' a outra --- ou o contrário, quando a
API/linguagem faz cópia.
\item \textbf{Confirme contra a assinatura real da função/método}, não
contra o que você acha que ela deveria fazer.
\end{enumerate}
\end{conceito}

\textbf{Exemplo resolvido (NumPy --- referência $\times$ cópia):}

\begin{center}
\begin{tabular}{@{}p{5.4cm}p{5.4cm}@{}}
\toprule
\textbf{Código} & \textbf{Execução mental} \\
\midrule
\texttt{a = np.array([1,2,3])} & \texttt{a} = [1,2,3] (original) \\
\texttt{b = a.view()} & \texttt{b} \textbf{compartilha} o buffer de \texttt{a}; \texttt{b.base is a} \\
\texttt{a[0] = 99} & muda o buffer compartilhado \\
\texttt{print(b[0])} & imprime \textbf{99} (reflete, pois é view) \\
\bottomrule
\end{tabular}
\end{center}

Se a segunda linha fosse \texttt{b = a.copy()}, \texttt{b[0]} continuaria
\textbf{1} depois de alterar \texttt{a[0]} --- porque \texttt{copy()} rompe o
vínculo com o buffer original. A pegadinha da FGV é exatamente trocar
\texttt{view} por \texttt{copy} no enunciado e esperar que você não perceba a
diferença de comportamento.

\begin{comosair}[Como treinar leitura ativa]
Toda vez que uma questão do quiz trouxer código, resista ao impulso de ler a
alternativa primeiro. Pegue papel, simule a execução, chegue a um valor, e só
depois compare com as alternativas --- isso elimina o efeito ``parece
certo''. Vale para Python, SQL, Java (Capítulo \ref{cap:java}) e
HTML/CSS/JS (Capítulo \ref{cap:frontend}): é a mesma disciplina em qualquer
linguagem.
\end{comosair}

\begin{jacaiu}
Spring/Spring Cloud/Spring Boot/Hibernate/JUnit (papéis); XML/XSLT/JSON;
Flutter (Dart); blockchain (o que não fica no bloco); SPA $\times$ PWA (ver
Capítulo \ref{cap:frontend}); NumPy \texttt{view} $\times$ \texttt{copy};
Kotlin/Swift; leitura de código pequeno em geral. Rode \texttt{./quiz.py
programacao}.
\end{jacaiu}

\section{Alta probabilidade / pesquisa extra}

\begin{itemize}
\item \textbf{Confluent Kafka} (mensageria/streaming): tópicos,
produtores/consumidores --- ver Capítulo \ref{cap:arquitetura}.
\item \textbf{API/Swagger (OpenAPI):} documenta e contrata REST; Swagger UI.
\item \textbf{RPA} (UiPath, Automation Anywhere): automação de tarefas de
interface.
\item \textbf{IA/LLM} entrou no edital --- ver Capítulo \ref{cap:atualidades}
para os fundamentos.
\end{itemize}

\begin{comosair}
Fixe uma frase por framework: Boot = configura e sobe sozinho (servidor
embarcado); Cloud = distribuído/service discovery; Hibernate =
objeto$\leftrightarrow$tabela; JUnit = testa; Spring = container/injeção.
Nativo mobile: Kotlin/Android, Swift/iOS (Objective-C e Java são os antigos).
Multiplataforma Dart = Flutter. Gatilhos: ``exclusivamente'', ``sempre'',
``estritamente equivalentes'', ``elimina a necessidade''. Em questão de
código, leia linha a linha --- a FGV muda um parâmetro só.
\end{comosair}
